\documentclass[11pt]{article}

\usepackage[margin=1in]{geometry}
\usepackage{amsmath,amssymb,amsthm}
\usepackage{mathtools}
\usepackage{booktabs}
\usepackage{bm}
\usepackage{enumitem}
\usepackage[round]{natbib}
\usepackage[hidelinks]{hyperref}

\newtheorem{proposition}{Proposition}
\newtheorem{lemma}{Lemma}
\newtheorem{assumption}{Assumption}
\newtheorem{remark}{Remark}

\DeclareMathOperator*{\argmin}{arg\,min}
\DeclareMathOperator*{\argmax}{arg\,max}
\newcommand{\Bin}{\mathrm{Bin}}
\newcommand{\E}{\mathbb{E}}
\newcommand{\Prob}{\mathbb{P}}
\newcommand{\one}{\mathbf{1}}
\newcommand{\Frechet}{\mathrm{Fr\acute{e}chet}}
\newcommand{\Nd}{N_{d}}
\newcommand{\Areas}{\mathcal{A}}

\title{\textbf{A Granular Eaton--Kortum Model of Spatial Sourcing\\[3pt]
\large Comparative Advantage at the Attraction-Area Level}}
\author{}
\date{\today}

\begin{document}

\maketitle

\begin{abstract}
In a continuum-of-varieties Eaton--Kortum economy every origin with a positive comparative
advantage supplies every buyer: the extensive margin of sourcing is degenerate. In the data most
region--sectors host no supplier at all --- inside the attraction areas that host at least one,
$91.5\%$ of employment zones have none. This paper makes the sourcing extensive margin a genuine
random object by giving each sector a \emph{finite} number $N_{s}$ of varieties while keeping a
continuum of firms within each variety and region, and by locating comparative advantage at the
level at which the data identify the trade-cost elasticity: the \emph{attraction area}. The two
departures are complementary. Area-level comparative advantage means that within an active area
the only thing separating an employment zone that hosts a supplier from one that does not is
distance; finiteness is what allows a zone with positive comparative advantage to host no supplier
at all. Together they deliver an exact reduced form --- the probability that a firm supplies is a
complementary log-log function of log distance and log own productivity with an
area$\times$sector fixed effect, with coefficients $\theta\alpha$ and $-\theta$ and \emph{no}
dependence on $N_{s}$ --- and a clean division of labour in identification: the distance gradient
identifies $\alpha$, the level of the zero count identifies $N_{s}$, and area-aggregated sourcing
shares identify comparative advantage.
\end{abstract}

\section{The degenerate extensive margin, and where the data identify $\alpha$}
\label{sec:motivation}

\subsection{The problem}

In the baseline model each upstream sector $s$ contains a continuum of varieties. Variety $\rho$
produced in region $l$ has a Fr\'echet productivity with shape $\theta$ and scale $T_{sl}$, and a
downstream buyer in region $r$ sources each variety from its cheapest origin. Ricardian selection
delivers the EK sourcing share
\begin{equation}
\gamma_{lrs}
=\frac{T_{sl}\,(w_{l}\tau_{lrs})^{-\theta}}{\Phi_{sr}},
\qquad
\Phi_{sr}\equiv\sum_{l'}T_{sl'}(w_{l'}\tau_{l'rs})^{-\theta},
\qquad
\tau_{lrs}=d_{lr}^{\alpha},
\label{eq:ek-share}
\end{equation}
with $w_{l}$ the upstream wage. Equation \eqref{eq:ek-share} is the probability that a
\emph{single} variety is won by origin $l$. With a continuum of varieties the law of large numbers
turns that probability into a deterministic realised share, so
\begin{equation}
\Prob\big(\text{$l$ supplies nothing to $r$ in $s$}\big)=\lim_{N\to\infty}(1-\gamma_{lrs})^{N}=0 .
\label{eq:degenerate}
\end{equation}
Every origin with $T_{sl}>0$ is active in every market. The observed zeros are not a model
outcome, and the extensive margin --- the object of interest --- has no content.

\subsection{Where the data identify the trade-cost elasticity}
\label{sec:fe}

The empirical extensive-margin regression carries a fixed effect at the \textbf{attraction area
$\times$ sector} level, an attraction area being the set of upstream employment zones (ZE) sharing
the same nearest downstream buyer. That choice determines what $\alpha$ is identified from.

\begin{lemma}[All-zero groups carry no information about the distance slope]
\label{lem:fe}
Index observations by firm $i$ within fixed-effect group $g=(a,s)$, with outcome $y_{i}$,
regressor $x_{i}$, weights $\omega_{i}>0$ and group effect $\phi_{g}$.
\begin{enumerate}[nosep,topsep=2pt,label=(\roman*)]
\item \emph{Complementary log-log.} The log-likelihood of a group in which $y\equiv0$ is
$\ell_{g}=-e^{\phi_{g}}\sum_{i}e^{\beta x_{i}}$, whose supremum over $\phi_{g}$ is attained as
$\phi_{g}\to-\infty$ with $\ell_{g}\to0$ and $\partial\ell_{g}/\partial\beta\to0$: the group drops
out of the concentrated likelihood and contributes \emph{exactly nothing} to $\hat\beta$.
\item \emph{Linear probability model.} By Frisch--Waugh--Lovell an all-zero group has
$\tilde y\equiv0$, contributing zero to the numerator and a positive amount to the denominator, so
$\hat\beta=\lambda\hat\beta^{+}$ with
$\lambda=\big(\sum_{a\in\Areas^{+}}\sum\omega\tilde x^{2}\big)/\big(\sum_{a}\sum\omega\tilde x^{2}\big)\in(0,1]$,
where $\Areas^{+}$ is the set of areas hosting at least one supplier.
\end{enumerate}
\end{lemma}

Under the cloglog link --- the link the model implies, Section~\ref{sec:reduced} --- $\alpha$ is
therefore identified \emph{exclusively} from variation \emph{within} attraction areas that host at
least one supplier. Under a linear probability model the empty areas enter only as a mechanical
dilution of the same within-area signal.

\subsection{The design}

Two departures, each answering one of the two problems.

\begin{enumerate}[nosep,topsep=2pt,label=(\alph*)]
\item \textbf{Granularity.} A finite number $N_{s}$ of varieties per sector makes a zero possible
at all.
\item \textbf{Area-level comparative advantage.} $T$ constant within an attraction area makes the
within-area extensive margin a pure trade-cost object, so what the model explains and what the
data identify are the same thing.
\end{enumerate}

\section{Model}
\label{sec:model}

The downstream block is unchanged: downstream regions $r=1,\dots,\Nd$ each hosting a
representative producer; nested-CES technology with labour/intermediate elasticity $\lambda$,
across-sector elasticity $\nu$ and within-sector elasticity $\nu_{s}$; demand elasticity
$\varepsilon$ and shifters $\delta_{r}$. Only the upstream supply side changes.

\subsection{Granular varieties, a continuum of firms}

\begin{assumption}[Granularity]
\label{as:granular}
Each upstream sector $s$ consists of $N_{s}\in\mathbb{N}$ horizontally differentiated varieties,
$\rho=1,\dots,N_{s}$. Within each variety $\rho$ and each ZE $l$ there is a \emph{continuum} of
firms whose productivities above $z$ form a Poisson process with mean $T_{sl}z^{-\theta}$. Draws
are independent across varieties, ZE and sectors.
\end{assumption}

The continuum of firms is what keeps the model tractable: the number of firms above $z$ is Poisson
with mean $T_{sl}z^{-\theta}$, so the \emph{most productive} firm of ZE $l$ for variety $\rho$ has
\begin{equation}
\Prob\big(\varphi_{l\rho s}\le z\big)=\Prob(\text{no firm above }z)=\exp\!\big(-T_{sl}z^{-\theta}\big),
\label{eq:champion}
\end{equation}
i.e.\ $\varphi_{l\rho s}\sim\Frechet(T_{sl},\theta)$ \emph{exactly}. One Fr\'echet draw per
(ZE, variety) \emph{is} the champion of a continuum of firms, not an approximation to it. Since
the iceberg cost scales all ZE-$l$ firms identically, that champion is the unique ZE-$l$ firm that
ever serves any buyer sourcing $\rho$ from $l$. \textbf{A firm is therefore a (ZE, variety)
champion}, with delivered cost $c_{l\rho s}(r)=w_{l}\tau_{lrs}/\varphi_{l\rho s}$, and the
probability that $l$ wins variety $\rho$ for buyer $r$ is exactly the EK share
\eqref{eq:ek-share}.

Only the number of \emph{varieties} is made finite; the firms are not discretised. This is the
``fixed number of independent Ricardian draws'' form of granular EK \citep[cf.][]{eks2012,bejk2003}.

\subsection{The two-tier geography}

\begin{assumption}[Area-level comparative advantage]
\label{as:geo}
Upstream ZE are partitioned into attraction areas, $\mathcal{L}=\bigsqcup_{a\in\Areas}\mathcal{L}_{a}$
with $a(l)$ the area of $l$, $\mathcal{L}_{a}$ being the ZE whose nearest downstream buyer is the
anchor of $a$. Comparative advantage is area-specific,
\begin{equation}
T_{sl}=T_{s,a(l)}\qquad\text{for every } l\in\mathcal{L}_{a},
\label{eq:T-area}
\end{equation}
with one parameter per (sector, area). Productivity draws remain \emph{independent across ZE}: two
ZE in the same area share the Fr\'echet \emph{scale}, not the realisation of their champions.
\end{assumption}

The last sentence matters. \eqref{eq:T-area} does not make ZE within an area interchangeable: each
runs its own Ricardian competition for each variety, and two ZE at different distances have
different win probabilities. What it removes is \emph{unobserved technological heterogeneity}
across ZE within an area. The sourcing share becomes
\begin{equation}
\gamma_{lrs}=\frac{T_{s,a(l)}\,(w_{l}\tau_{lrs})^{-\theta}}{\Phi_{sr}},
\qquad
\Phi_{sr}=\sum_{a'}T_{s,a'}\sum_{l'\in\mathcal{L}_{a'}}(w_{l'}\tau_{l'rs})^{-\theta}.
\label{eq:share-area}
\end{equation}

Two ZE $l,l'$ in the same area and the same buyer therefore satisfy
\begin{equation}
\ln\frac{\gamma_{lrs}}{\gamma_{l'rs}}=-\theta\alpha\big(\ln d_{lr}-\ln d_{l'r}\big)-\theta\ln\frac{w_{l}}{w_{l'}},
\label{eq:within}
\end{equation}
the technology term and the price index having cancelled. Under the wage normalisation $w\equiv1$
the within-area sourcing profile is a function of $(\theta,\alpha)$ and observed distances alone.

\subsection{The extensive margin}

Across the $N_{s}$ varieties the win events are i.i.d., so with $K_{lrs}$ the number of sector-$s$
varieties buyer $r$ sources from origin $l$,
\begin{equation}
K_{lrs}\sim\Bin\big(N_{s},\gamma_{lrs}\big),
\qquad
\Prob(K_{lrs}=0)=(1-\gamma_{lrs})^{N_{s}}>0 .
\label{eq:binomial}
\end{equation}
Unlike \eqref{eq:degenerate} the zero-supplier probability is strictly positive and endogenous: it
falls with comparative advantage, rises with distance through $\tau_{lrs}=d_{lr}^{\alpha}$, and
falls with $N_{s}$. Letting $q_{ls}$ be the probability that $l$ wins a given variety
\emph{somewhere} in the downstream industry (the union over buyers), the industry-wide count is
$K_{ls}\sim\Bin(N_{s},q_{ls})$.

\begin{remark}[Dependence across buyers]
\label{rem:dependence}
For a given variety the champion costs are common across buyers; only $\tau_{lrs}$ differs. Win
events are therefore positively dependent across buyers. The per-market marginal
\eqref{eq:binomial} is nonetheless exact, and the union probability $q_{ls}$ is evaluated by
simulation.
\end{remark}

\subsection{A hierarchy of zeros}
\label{sec:hierarchy}

Assumption~\ref{as:geo} sorts observed zeros into two economically distinct classes. This is the
cleanest statement of what the model does and does not endogenise.

\begin{table}[htbp]
\centering
\small
\begin{tabular}{@{}p{0.20\textwidth}p{0.24\textwidth}p{0.22\textwidth}p{0.26\textwidth}@{}}
\toprule
\textbf{Zero} & \textbf{Treatment} & \textbf{Probability} & \textbf{Role}\\
\midrule
Whole attraction area $a$ has no supplier in $s$
& \emph{Structural}: $T_{s,a}=0$, $(s,a)$ outside the active set
& $1$ by construction
& None --- absorbed by the $a\times s$ fixed effect (Lemma~\ref{lem:fe}); no parameter estimated\\[4pt]
ZE $l$ inside an active area has no supplier
& \emph{Endogenous}: inherits $T_{s,a}>0$, own distance $d_{l}$
& $(1-q_{ls})^{N_{s}}$
& \textbf{The identifying variation}: disciplined by distance alone \eqref{eq:within}\\
\bottomrule
\end{tabular}
\caption{The two classes of extensive-margin zeros. The second row is exactly the set of
employment zones that host firms but no supplier. Under ZE-level comparative advantage they had to
be assigned $T\equiv0$, making their zeros exogenous; under Assumption~\ref{as:geo} they inherit
their area's comparative advantage and their zeros become model predictions.}
\label{tab:hierarchy}
\end{table}

\section{What the model implies for the data}
\label{sec:reduced}

\subsection{The firm-level reduced form}

The empirical extensive-margin regression is run at the firm level, with the firm's own
productivity as a control. The model delivers that specification exactly.

\begin{proposition}[Firm-level reduced form]
\label{prop:firm}
Let a firm be a (ZE, variety) champion with productivity $z$, and let $a$ be its area's downstream
anchor. Writing $\Phi^{-l}_{sa}=\sum_{l'\neq l}T_{sl'}(w_{l'}\tau_{l'as})^{-\theta}$ for the
competitors' cost scale,
\begin{equation}
\Prob\big(\text{the firm supplies}\mid z\big)=\exp\!\Big(-\Phi^{-l}_{sa}\big(w_{l}\tau_{las}\big)^{\theta}z^{-\theta}\Big),
\label{eq:firm-win}
\end{equation}
so the probability that it does \emph{not} supply is a complementary log-log function
\begin{equation}
\Prob\big(\text{no supply}\mid z\big)=1-\exp\!\big(-\exp(\eta)\big),
\qquad
\eta=\underbrace{\ln\Phi^{-l}_{sa}}_{\text{area}\times\text{sector FE}}
+\theta\ln w_{l}
+\boxed{\theta\alpha}\,\ln d_{la}
+\boxed{(-\theta)}\,\ln z .
\label{eq:firm-cloglog}
\end{equation}
The variety count $N_{s}$ does not appear.
\end{proposition}

Three consequences, and they are the backbone of the estimation.

\begin{enumerate}[nosep,topsep=3pt,label=(\alph*)]
\item \textbf{The empirical specification is the model, not an approximation to it.} A cloglog on
the ``does not supply'' indicator, on log distance and log own productivity, with an
area$\times$sector fixed effect, is exactly \eqref{eq:firm-cloglog}. The distance coefficient is
$+\theta\alpha$, so $\alpha=\hat\beta_{\ln d}/\theta$ with no attenuation correction.
\item \textbf{$\alpha$ and $N_{s}$ separate exactly.} $N_{s}$ is absent from
\eqref{eq:firm-cloglog}: conditioning on the firm removes it. The regression identifies $\alpha$;
the count level identifies $N_{s}$; there is no residual coupling.
\item \textbf{The productivity coefficient is a free over-identifying test.} It equals $-\theta$.
Extending the moment vector to include it would let $\theta$ be estimated rather than calibrated.
\end{enumerate}

Table~\ref{tab:convention} confirms the proposition numerically and shows that the outcome
convention is dictated by the unit of observation rather than chosen: the reversed outcome
recovers neither coefficient.

\begin{table}[htbp]
\centering
\small
\begin{tabular}{@{}lrr@{}}
\toprule
outcome & $\hat\beta_{\ln d}$ \ (truth $\theta\alpha=0.30$) & $\hat\beta_{\ln z}$ \ (truth $-\theta=-1.0$)\\
\midrule
does not supply & $+0.327$ & $-1.051$\\
supplies        & $-0.426$ & $+1.452$\\
\bottomrule
\end{tabular}
\caption{Simulation of the exact object, $\theta=1$, $\alpha=0.30$, firms as (ZE $\times$ variety)
champions.}
\label{tab:convention}
\end{table}

\subsection{The count moment}
\label{sec:count}

The zeros themselves are a cell-level object. Let $\mathcal{L}^{+}_{s}$ be the ZE inside active
attraction areas. The empirical moment is the share of those ZE hosting at most $n$ suppliers,
\begin{equation}
\bar G_{s}(n)=\frac{1}{|\mathcal{L}^{+}_{s}|}\#\big\{l\in\mathcal{L}^{+}_{s}:K_{ls}\le n\big\},
\qquad
\E\big[\bar G_{s}(0)\big]=\frac{1}{|\mathcal{L}^{+}_{s}|}\sum_{l}(1-q_{ls})^{N_{s}},
\label{eq:count}
\end{equation}
the count analogue of the number-of-buyers-per-seller moment of \citet{lmm2023}. It is strictly
decreasing in $N_{s}$, which is what identifies the variety count: given the sourcing shares ---
pinned by area-aggregated value shares and by $\alpha$, neither of which involves $N_{s}$ --- the
only free object shaping \eqref{eq:count} is $N_{s}$.

The support is restricted to active areas for the same reason the fixed effect discards the
others: ZE in empty areas are structural zeros carrying no information about $(\alpha,N_{s})$.

\subsection{Bounds on the variety count}
\label{sec:bounds}

Because the iceberg cost never changes which firm is cheapest within a ZE, each variety generates
one distinct supplying firm per winning origin. So if $N^{\mathrm{obs}}_{s}$ is the number of
distinct upstream firms observed to supply the downstream industry,
$N^{\mathrm{obs}}_{s}=\sum_{\rho}\#\{\text{origins winning }\rho\text{ somewhere}\}\in[N_{s},N_{s}\Nd]$,
hence
\begin{equation}
\frac{N^{\mathrm{obs}}_{s}}{\Nd}\;\le\;N_{s}\;\le\;N^{\mathrm{obs}}_{s} .
\label{eq:bounds}
\end{equation}
The bounds are tight: the upper when every variety is globally sourced from a single origin, the
lower when each of the $\Nd$ buyers sources it from a different origin. Since
$\E[\sum_{l}K_{ls}]=N_{s}\sum_{l}q_{ls}$ with $\sum_{l}q_{ls}\in[1,\min(\Nd,R)]$, matching the
observed firm count gives a second, closed-form estimate
$N^{\mathrm{count}}_{s}=N^{\mathrm{obs}}_{s}/\sum_{l}q_{ls}$ that lies in \eqref{eq:bounds} by
construction --- a free over-identifying check.

\subsection{Identification}
\label{sec:identification}

\begin{table}[htbp]
\centering
\small
\begin{tabular}{@{}llp{0.36\textwidth}@{}}
\toprule
\textbf{Parameter} & \textbf{Identifying variation} & \textbf{Moment}\\
\midrule
$\alpha$ & \emph{Within}-area, across-firm gradient of supplying in distance
& firm-level cloglog distance coefficient \eqref{eq:firm-cloglog}\\
$T_{s,a}$, $a\in\Areas^{+}_{s}$ & \emph{Across}-area value shares
& area-aggregated sourcing shares $\gamma_{s,a}=\sum_{l\in\mathcal{L}_{a}}\gamma_{ls}$\\
$N_{s}$ & \emph{Level} of the empty-ZE share inside active areas
& count moment $\bar G_{s}(0)$ \eqref{eq:count}\\
$\theta$ & Gradient of supplying in own productivity
& $\hat\beta_{\ln z}=-\theta$ --- currently calibrated, hence a test\\
$\Areas^{+}_{s}$ & Which areas host any supplier
& \emph{Not modelled} --- taken as given, matching the fixed effect's information set\\
\bottomrule
\end{tabular}
\caption{The channels are near-orthogonal by construction: $\alpha$ reads a within-area slope that
$T$ cannot touch because it is constant within the area; $T$ reads across-area aggregates;
$N_{s}$ reads a level that the fixed effect strips out of the $\alpha$ regression.}
\label{tab:identification}
\end{table}

\subsection{Degrees of freedom}

Fix a sector and let $A^{+}=|\Areas^{+}_{s}|$ and $R^{\mathrm{act}}_{a}$ be the number of ZE in
area $a$ with a strictly positive observed share. The $T$ block has $A^{+}-1$ free elements after
the reference normalisation, against $\sum_{a}R^{\mathrm{act}}_{a}$ share moments, so
Assumption~\ref{as:geo} converts what were free parameters into
\begin{equation}
\sum_{a\in\Areas^{+}}\big(R^{\mathrm{act}}_{a}-1\big)
\quad\text{over-identifying restrictions,}
\end{equation}
one per active ZE beyond the first in each area. Each is a within-area statement: given the area
aggregate, the split across its ZE must follow the distance profile \eqref{eq:within}. Under
ZE-level comparative advantage this count is zero.

\section{Estimation}
\label{sec:estimation}

\subsection{Moments and simulation}

The moment vector has six blocks: the aggregate labour share, industry shares $\pi_{s}$, regional
sales shares $\pi_{r}$, the firm-level distance coefficients of \eqref{eq:firm-cloglog}, the
area-aggregated sourcing shares, and the count moment $\bar G_{s}(0)$. The economy is simulated
with exactly $N_{s}$ varieties per sector, replicated, and moments are computed on the realised
economies and averaged.

Two properties keep this tractable. First, the winner of a variety is the argmin over ZE of that
variety's own draws, so it does not depend on how many varieties exist: one Ricardian solve at
pool width serves every candidate $N_{s}$ as a prefix, with no re-drawing. Second:

\begin{lemma}[$q_{ls}$ is free of $N_{s}$]
\label{lem:qfree}
Winning a variety is the comparison
$w_{l}\tau_{lrs}/\varphi_{l\rho s}<\min_{l'\neq l}w_{l'}\tau_{l'rs}/\varphi_{l'\rho s}$, so
$q_{ls}=q_{ls}(T_{s\cdot},w,\tau,\theta)$ and nothing else. The variety count, the price index,
downstream costs and expenditures do not appear.
\end{lemma}

The economics is that winning is a statement about \emph{relative costs}, whereas $N_{s}$ and
$P_{sr}$ are aggregators built after the winners are known: expenditure allocates value across
winners, it does not choose them. The one condition that would break the lemma is endogenous
upstream wages, since $w$ enters the comparison directly; wages are taken as data.

\subsection{The variety count}

$N_{s}$ is estimated by the same weighted objective as every other parameter, concentrated rather
than searched jointly:
\begin{equation}
\widehat N(\theta_{c})=\argmin_{N\in\prod_{s}[N^{\mathrm{lo}}_{s},N^{\mathrm{hi}}_{s}]\cap\mathbb{N}}
r(\theta_{c},N)'W\,r(\theta_{c},N),
\label{eq:profile}
\end{equation}
with the bounds of \eqref{eq:bounds}. This is the profiled extremum estimator: the result is
identical to joint minimisation over $(\theta_{c},N)$, $\widehat N_{s}$ is an integer at every
evaluation so integrality is never relaxed and never rounded, and by Lemma~\ref{lem:qfree} the
inner minimisation needs no re-simulation. Since $\bar G_{s}(0;n)$ is closed form and strictly
decreasing in $n$, \eqref{eq:profile} reduces to a monotone integer bisection.

Clamping at either bound is informative rather than benign: the upper bound binding means the model
cannot generate enough sparsity even when every variety is sourced from a single origin.

\subsection{The price index}
\label{sec:priceindex}

The one place where granularity meets the value block requires an explicit assumption, and it is
worth stating with its magnitude.

With $N_{s}$ varieties the within-sector price index is a finite CES aggregate,
$P_{sr}^{1-\nu_{s}}=\frac{1}{N_{s}}\sum_{\rho}p_{\rho rs}^{1-\nu_{s}}$, whose expectation is
invariant to $N_{s}$ --- the love-of-variety level is absorbed into the within-sector scale of $T$,
which is a free gauge (Appendix~\ref{app:price}). But $P_{sr}$ is a \emph{convex} transform of that
average, so the realised index exceeds its continuum limit by Jensen, and the gap closes only as
fast as the average concentrates. At the calibration $\theta=1$, $\nu_{s}=1.5$ the regularity
condition $\theta>2(\nu_{s}-1)$ reads $1>1$: it fails at equality, $(c^{*})^{1-\nu_{s}}$ has
infinite variance, and the average converges at $\sqrt{N/\log N}$ rather than $\sqrt{N}$.

\begin{table}[htbp]
\centering
\small
\begin{tabular}{@{}lrrrrrrr@{}}
\toprule
$N_{s}$ & $25$ & $50$ & $100$ & $200$ & $400$ & $1000$ & $4000$\\
\midrule
$\E[\widehat P]$ vs continuum & $+18.6\%$ & $+11.1\%$ & $+6.5\%$ & $+3.8\%$ & $+2.2\%$ & $+1.0\%$ & $+0.3\%$\\
$\mathrm{sd}(\widehat P)/P$ & $0.48$ & $0.35$ & $0.27$ & $0.20$ & $0.15$ & $0.11$ & $0.06$\\
\bottomrule
\end{tabular}
\caption{The realised granular price index against its continuum limit, $\theta=1$, $\nu_{s}=1.5$.}
\label{tab:price-jensen}
\end{table}

So $N_{s}$ does reach the value block, and at $N_{s}\approx10^{2}$ by about $6\%$. The estimator
adopts the \textbf{certainty-equivalent} reading: downstream firms optimise against
$\E[P^{1-\nu_{s}}]^{1/(1-\nu_{s})}$, neglecting granular price randomness --- the
general-equilibrium embedding of \citet{lmm2023}. Two things make this the right default. Most of
the $6\%$ would not have been identifying variation: a price-level shift is observationally
equivalent to a technology shift (Appendix~\ref{app:price}), so the sector technology terms absorb
it and only the cross-sector pattern of $N_{s}$ survives as signal. And by Lemma~\ref{lem:qfree}
none of it touches $q_{ls}$ or the extensive margin, so the identification of $\alpha$ and $N_{s}$
is unaffected either way.

What the assumption must not do is survive into counterfactuals. The love-of-variety propagation
channel lives in the realised index --- a shock that removes suppliers raises $P_{sr}$ downstream
through exactly that channel, and it is absent from the continuum model. Any counterfactual moving
the variety set must use the realised index.

\section{Key quantitative facts}
\label{sec:facts}

\subsection{The extensive margin is a within-area object; the intensive margin is not}

Inside attraction areas hosting at least one supplier there are $1161$ ZE--sector cells, of which
$99$ host a supplier: \textbf{$91.5\%$ are extensive-margin zeros}, and every one of those cells
contains at least one firm. Entering the fixed effect as a regressor,
area$\times$sector dummies explain
\begin{itemize}[nosep,topsep=2pt]
\item $5.9\%$ of the variance of the sourcing share in \emph{levels} --- a variable that is
$91.5\%$ zeros, hence dominated by the extensive margin, so \textbf{the extensive margin is
$94\%$ within-area};
\item $94.4\%$ of the variance of \emph{log} shares among active cells --- so \textbf{the intensive
margin is $94\%$ between-area}.
\end{itemize}
The data split along exactly the seam Assumption~\ref{as:geo} cuts: $T_{s,a}$ for the intensive
margin, distance and granularity for the extensive margin.

\subsection{Size dominates within-area variation}

Regressing observed sourcing shares within area$\times$sector on log distance and log firm count:

\begin{table}[htbp]
\centering
\small
\begin{tabular}{@{}lrrrr@{}}
\toprule
& \multicolumn{2}{c}{share, levels} & \multicolumn{2}{c}{log share, active only}\\
\cmidrule(lr){2-3}\cmidrule(lr){4-5}
$\ln d$ & $-0.017^{***}$ & $-0.011^{***}$ & $-0.284$ & $-0.130$\\
        & $(0.004)$ & $(0.002)$ & $(0.197)$ & $(0.136)$\\
$\ln(\text{firm count})$ & & $0.010^{***}$ & & $0.492^{***}$\\
        & & $(0.003)$ & & $(0.126)$\\
\midrule
within $R^{2}$ & $0.108$ & $0.225$ & $0.084$ & $0.330$\\
Observations & $1161$ & $1161$ & $99$ & $99$\\
\bottomrule
\end{tabular}
\caption{Within-area diagnostic. Standard errors clustered by area$\times$sector.}
\label{tab:diagnostic}
\end{table}

Size carries four times the explanatory power of distance and, when added, moves the distance
coefficient toward zero in both specifications: within an area, closer ZE are \emph{bigger}, so
omitting size makes distance look more powerful than it is. Both coefficients are attenuated by
the same factor --- conditional on being active, a ZE's observed share is quantised toward $1/N_{s}$
--- so their \emph{ratio} is the interpretable object, and only in the level specification, where
$\E[\hat\gamma]=\gamma$ holds unconditionally. With $\theta=1$ it gives
$\alpha/\eta_{\text{size}}\approx1.1$.

The implication for the model is direct: any within-area heterogeneity other than distance must
enter explicitly, either as observed covariates scaling $T$ with a common loading or as a random
deviation integrated out; it must never become a free parameter per ZE, which would restore exact
just-identification and remove the over-identifying content of Assumption~\ref{as:geo}.

\section{Extensions}

\paragraph{ZE-level deviations.} Assumption~\ref{as:geo} can be relaxed to
$T_{sl}=T_{s,a(l)}\nu_{sl}$ with $\E[\nu]=1$ and $\mathrm{Var}(\nu)=\sigma^{2}_{s}$, calibrated on
the dispersion of the ZE-level shares. With $\nu$ Gamma the supplier count becomes negative
binomial and the whole count family stays closed form. The deviation must be a \emph{distribution},
never a parameter vector: a free $\nu$ per ZE is undefined precisely at the empty ZE the model
exists to explain, and setting it to one there is biased downward by selection,
$\E[\nu\mid K=0]=(1+\sigma^{2}\mu)^{-1}$.

\paragraph{Firm counts.} Proposition~\ref{prop:firm} identifies each observed firm with a
(ZE, variety) champion, which gives every cell exactly $N_{s}$ firms while observed counts vary by
orders of magnitude. Closing that gap means giving the model a firm count --- a productivity floor
$\bar z_{s}$ makes the count in $(l,s,\rho)$ Poisson with mean $T_{sl}\bar z_{s}^{-\theta}$, so it
varies with $T_{sl}$ as in the data, and the observed count becomes a new and strong moment for
$T$. The $\ln z$ coefficient of Proposition~\ref{prop:firm}(c) is the cheapest diagnostic of
whether this is needed.

\appendix

\section{Proofs}
\label{app:proofs}

\subsection*{Lemma~\ref{lem:fe}}

(i) With $\Prob(y=1)=1-\exp(-\exp(\phi_{g}+\beta x))$, a group in which $y\equiv0$ contributes
$\ell_{g}=\sum_{i}\ln\exp(-e^{\phi_{g}+\beta x_{i}})=-e^{\phi_{g}}\sum_{i}e^{\beta x_{i}}$. This is
increasing in $-\phi_{g}$ with supremum $0$ as $\phi_{g}\to-\infty$, at which
$\partial\ell_{g}/\partial\beta=-e^{\phi_{g}}\sum_{i}x_{i}e^{\beta x_{i}}\to0$. The group therefore
drops from the concentrated likelihood.

(ii) By Frisch--Waugh--Lovell,
$\hat\beta=\sum_{g,i}\omega\tilde x\tilde y/\sum_{g,i}\omega\tilde x^{2}$ with $\tilde\cdot$ the
weighted within-group demeaning. A group with $y\equiv0$ has $\tilde y\equiv0$, so it contributes
zero to the numerator and $\sum\omega\tilde x^{2}>0$ to the denominator. Splitting the denominator
between $\Areas^{+}$ and its complement gives $\hat\beta=\lambda\hat\beta^{+}$. \hfill$\square$

\subsection*{Proposition~\ref{prop:firm}}

Condition on the firm's productivity $z$; its delivered cost to buyer $a$ is $c=w_{l}\tau_{las}/z$.
Each competitor $l'$ has champion cost $w_{l'}\tau_{l'as}/\varphi_{l'}$ with
$\varphi_{l'}\sim\Frechet(T_{sl'},\theta)$, so
$\Prob(w_{l'}\tau_{l'as}/\varphi_{l'}>x)=\exp(-T_{sl'}(w_{l'}\tau_{l'as})^{-\theta}x^{\theta})$.
By independence across competitors the minimum competing cost $M$ satisfies
$\Prob(M>x)=\exp(-\Phi^{-l}_{sa}x^{\theta})$. The firm supplies iff $M>c$, giving
\eqref{eq:firm-win}. Taking the complement,
$\Prob(\text{no supply}\mid z)=1-\exp(-\exp(\eta))$ with
$\eta=\ln\Phi^{-l}_{sa}+\theta\ln(w_{l}\tau_{las})-\theta\ln z$, and $\tau_{las}=d_{la}^{\alpha}$
gives \eqref{eq:firm-cloglog}. No step involves the variety count. \hfill$\square$

\subsection*{Lemma~\ref{lem:qfree}}

Origin $l$ wins variety $\rho$ for buyer $r$ iff
$w_{l}\tau_{lrs}/\varphi_{l\rho s}<\min_{l'\neq l}w_{l'}\tau_{l'rs}/\varphi_{l'\rho s}$, an event
measurable with respect to $\{\varphi_{\cdot\rho s}\}$, $w$ and $\tau$ alone; the $\varphi$ have
distribution governed by $(T_{s\cdot},\theta)$. The win-somewhere event is the union over $r$ of
those events, hence measurable with respect to the same objects. Neither the variety count nor any
price or expenditure aggregate enters. \hfill$\square$

\subsection*{Equation \eqref{eq:champion}}

The number of firms in $(l,s,\rho)$ with productivity above $z$ is Poisson with mean
$T_{sl}z^{-\theta}$, so the probability that none exceeds $z$ --- equivalently that the maximum is
below $z$ --- is $\exp(-T_{sl}z^{-\theta})$, the $\Frechet(T_{sl},\theta)$ CDF. \hfill$\square$

\section{The price index, love of variety, and the normalisation}
\label{app:price}

\subsection*{The minimum delivered cost}

For variety $\rho$ in market $(r,s)$, origin $l$'s champion cost satisfies
$\Prob(c_{l\rho s}(r)>x)=\exp(-T_{sl}(w_{l}\tau_{lrs})^{-\theta}x^{\theta})$, so by independence
across origins the minimum $c^{*}_{\rho}$ is Weibull,
\begin{equation}
\Prob\big(c^{*}_{\rho}>x\big)=\exp\!\big(-\Phi_{sr}x^{\theta}\big),
\qquad
\E\big[(c^{*}_{\rho})^{k}\big]=\Phi_{sr}^{-k/\theta}\,\Gamma\!\Big(1+\tfrac{k}{\theta}\Big),
\quad k>-\theta .
\label{eq:app-weibull}
\end{equation}

\subsection*{Love of variety and its normalisation}

With the CES aggregate taken as a sum, $\E[P_{sr}^{1-\nu_{s}}]=N_{s}\Gamma(1+\tfrac{1-\nu_{s}}{\theta})\mu^{\nu_{s}-1}\Phi_{sr}^{-\frac{1-\nu_{s}}{\theta}}$,
so the certainty-equivalent index $\widetilde P_{sr}=(\E[P_{sr}^{1-\nu_{s}}])^{1/(1-\nu_{s})}$
satisfies
\begin{equation}
\widetilde P_{sr}=N_{s}^{\frac{1}{1-\nu_{s}}}\times\underbrace{\Gamma\!\Big(1+\tfrac{1-\nu_{s}}{\theta}\Big)^{\frac{1}{1-\nu_{s}}}\mu^{-1}\Phi_{sr}^{-1/\theta}}_{N_{s}\text{-free}} .
\label{eq:loveofvariety}
\end{equation}
Since $1-\nu_{s}<0$ the factor is decreasing in $N_{s}$: more varieties lower the sector price
index. Rescaling $\widetilde T_{s,a}=T_{s,a}N_{s}^{\kappa}$ with $\kappa=\theta/(1-\nu_{s})$ maps
$\Phi_{sr}\mapsto N_{s}^{\kappa}\Phi_{sr}$ and cancels the two powers exactly, leaving
$\widetilde P_{sr}$ invariant to $N_{s}$. Equivalently, taking the CES aggregate as an
\emph{average} rather than a sum implements the same normalisation.

The normalisation is consistent because it acts on a gauge the sourcing shares do not see: the
share \eqref{eq:share-area} is a ratio of $T$'s within a sector, so the extensive margin, the count
moment and the within-area contrast \eqref{eq:within} are all unchanged. It is also a gauge the
estimator already fixes, by normalising one reference area per sector.

\subsection*{Why the removed term is nonetheless economic}

The elasticity $\partial\ln\widetilde P_{sr}/\partial\ln N_{s}=1/(1-\nu_{s})<0$ is a genuine gain
from variety. Three implications. First, the cross-section cannot separate $N_{s}$ from technology
using prices alone --- a sector that is cheap because it has many varieties looks identical, in
expenditure data, to one that is cheap because it is productive --- which is why the count moment
is indispensable. Second, the normalisation fixes a \emph{level} at the calibration point and does
not alter the counterfactual \emph{response} of prices to a change in the variety set, provided one
maps back to the physical $T$ before running such a counterfactual. Third, it is a live propagation
channel: a shock that removes suppliers in one region raises $P_{sr}$ downstream through exactly
this elasticity, and that channel is absent from the continuum model.

\section{The count moment in estimation}
\label{app:count}

\paragraph{Which thresholds.} Writing $\mu_{ls}=N_{s}q_{ls}$ and using the Poisson approximation,
$\partial\bar G_{s}(K)/\partial N_{s}=-|\mathcal{L}^{+}_{s}|^{-1}\sum_{l}q_{ls}e^{-\mu_{ls}}\mu_{ls}^{K}/K!$,
which peaks at $K\approx\mu_{ls}$ and vanishes as $\bar G_{s}(K)\to1$. Saturated thresholds carry
no identifying variation and a small, poorly-estimated variance, so they must be trimmed. The
empty-ZE share $\bar G_{s}(0)$ is the most interpretable and typically near-optimal threshold: it
\emph{is} the extensive margin the model was built to explain.

\paragraph{Bootstrap design.} The moment is a proportion over the ZE of active attraction areas,
empty ZE included, since those are the $K=0$ mass. Two errors to avoid, in opposite directions:
restricting to ZE that host a supplier computes a share \emph{among} active ZE, whose variance is
not that of the extensive margin; and including ZE from empty areas inflates the denominator with
mechanically-explained zeros that the fixed effect has already discarded. The resampling unit is
the supplier firm, pooled over the geography of active areas, with the denominator held fixed --- a
ZE with one or two firms can then draw zero and flip into the empty set, which is precisely the
required variance. Because the same firms generate both the sourcing shares and the counts, a
single joint bootstrap is needed to estimate the cross-covariance blocks that the efficient weight
matrix requires.

\begin{thebibliography}{9}

\bibitem[Bernard et al.(2003)]{bejk2003}
Bernard, A., J. Eaton, J.~B. Jensen, and S. Kortum (2003).
``Plants and Productivity in International Trade.'' \emph{American Economic Review} 93(4), 1268--1290.

\bibitem[Eaton et al.(2012)]{eks2012}
Eaton, J., S. Kortum, and S. Sotelo (2012).
``International Trade: Linking Micro and Macro.'' NBER Working Paper 17864.

\bibitem[Lenoir et al.(2023)]{lmm2023}
Lenoir, C., J. Martin, and I. M\'ejean (2023).
``Search Frictions in International Goods Markets.'' \emph{Journal of the European Economic
Association} 21(1), 326--366.

\end{thebibliography}

\end{document}
